\section{Axioms and Semantics}\label{sec:mathematical-framework}

\subsection{Fibrational Semantics}

The path-object and fibrational part of this semantic layer is adapted from
Helfer's first-order homotopical logic
\cite[Sections~2--7]{helfer2019fol}. In particular, we use its replacement of
diagonal equality by a path-space fibration, its Grothendieck-fibrational
formulation of first-order semantics, and its homotopy-invariance theorem.
We do not reproduce Helfer's full development of suitable model categories,
Kan complexes, \(h^{=}\)-fibrations, or two-dimensional fibrations. Instead, we
extract the semantic obligations consumed by HCP and add the project-specific
public-classifier and cryptographic layers below. Helfer's external invariance
theorem is cited as background rather than reproved in Lean.
Lean does formalize the typed dependency needed to apply that result: a
multi-sorted algebraic signature, finite-product interpretations, homotopy
homomorphisms coherent with every function symbol, and componentwise
homotopy equivalence between cofibrant-fibrant interpretations. The
closed-formula consequence of Helfer's theorem remains an explicit hypothesis,
from which the HCP classifier bridge is proved in
\cref{prop:helfer-hcp-bridge}.
We use the usual fibrational notation
\[
  p:E\longrightarrow B,
  \qquad
  E_b=p^{-1}(b)
\]
for a total space of hidden data over a base of public or contextual data.
When \(x,y\in E_b\), equality in the semantic layer is interpreted through a
path object rather than by strict equality in the ambient presentation. We
write \(\mathcal P_X(x,y)\) for the path object from \(x\) to \(y\) in a space
or type \(X\), and set
\[
  x\simeq_b y
  \quad\Longleftrightarrow\quad
  \mathcal P_{E_b}(x,y)\neq\varnothing .
\]
For a dependent family \(F:E\to\mathsf{Type}\), a path
\(\gamma\in\mathcal P_E(x,y)\) supplies the transport map
\[
  \gamma_{\ast}:F_x\longrightarrow F_y.
\]
Thus the primitive semantic data are fibers, paths, and transport along paths.
The cryptographic objects introduced later live over this interface; they do
not replace it with byte-level equality.

We also use the corresponding categorical fibration vocabulary. For a functor
\(p:E\to B\), a morphism \(f:U\to V\) in \(E\) is \(p\)-Cartesian when every
arrow \(\psi:W\to V\) and base arrow \(u:p(W)\to p(U)\) satisfying
\[
  p(\psi)=p(f)\circ u
\]
admit a unique factor \(\widetilde\psi:W\to U\) with
\[
  p(\widetilde\psi)=u,
  \qquad
  f\circ\widetilde\psi=\psi.
\]
A Grothendieck fibration is then a functor
with Cartesian lifts for every base arrow. This is the categorical form of the
transport principle used by the hidden-representative semantics. Helfer's full
first-order interface imposes additional fiberwise products, coproducts,
exponentials, stable indexed sums and products, and equality objects
\cite[Definition~4.3]{helfer2019fol}; the definition here records only the
underlying fibration vocabulary needed for the local construction.

\begin{proposition}[Fibration copies preserve the categorical fibration surface]\label{prop:fibration-copy-equivalence}
  Let \(p:E\to B\) be a Grothendieck fibration. For each tag \(n\), let
  \(E^{(n)}\) be a tagged copy of \(E\) and let
  \(p^{(n)}:E^{(n)}\to B\) be the induced fibration. Then
  \(p^{(n)}\) is again a Grothendieck fibration, and \(p\) and \(p^{(n)}\) are
  equivalent as total categories through functors commuting strictly with the
  projections to \(B\).
\end{proposition}

\begin{proof}
  The tagged copy has the same objects and morphisms as \(E\), with only the
  external tag changed. The forgetful functor \(E^{(n)}\to E\) and the tagging
  functor \(E\to E^{(n)}\) compose to the identity up to the evident natural
  isomorphisms. Cartesian lifts in \(E^{(n)}\) are transported directly from
  Cartesian lifts in \(E\), so the induced functor remains a Grothendieck
  fibration. The two functors give the stated categorical equivalence over the
  base \(B\). This tagged-copy equivalence is only a relabelling construction;
  it is not Helfer's model-categorical homotopy equivalence of algebraic
  interpretations and does not by itself imply hidden-kernel ambiguity.
\end{proof}

\subsection{Path-Object Semantics for Equality}

Following Helfer's homotopical interpretation of intuitionistic first-order
logic \cite[Section~2 and Definitions~6.3 and~6.6]{helfer2019fol}, let \(A\)
be a Kan complex, or more generally an object in a suitable model category,
and choose a path-object factorization
\[
  \Delta_A:A\overset{\sim}{\longrightarrow}PA
  \twoheadrightarrow A\times A.
\]
In the simplicial realization one may take \(PA=A^I\), with the second map
given by endpoint evaluation. For points \(x,y\in A\), write
\[
  \mathcal P_A(x,y)
  = PA\times_{A\times A}\{(x,y)\}.
\]
The interpretation of the atomic equality formula \(x=y\) is this pullback.
Its semantic inhabitation condition is therefore
\[
  x\simeq_A y
  \quad\Longleftrightarrow\quad
  \mathcal P_A(x,y)\neq\varnothing .
\]
As an elementary local consequence, if the chosen path data also carry
constant paths, reversal, and composition, then \(\simeq_A\) is reflexive,
symmetric, and transitive. This equivalence-relation observation is not being
quoted as a separate theorem of Helfer.

Helfer's homotopy-invariance theorem applies under its stated categorical
hypotheses. Let \(M,N:\sigma\to\mathcal C_{\mathrm f}\) be base
interpretations of an algebraic signature in a suitable model category, let
\(\widehat M,\widehat N\) be homotopical interpretations over them, and let
\(M\simeq N\) be a homotopy equivalence of the base interpretations. Then every
formula-in-context has homotopy-equivalent interpretation objects over the
induced context equivalence. In particular,
\[
  M\simeq N
  \quad\Longrightarrow\quad
  \bigl(M\models\varphi \Longleftrightarrow N\models\varphi\bigr)
\]
for every closed first-order formula \(\varphi\)
\cite[Definition~7.1 and Theorem~12.3]{helfer2019fol}. We use this external
theorem as semantic background; the local Lean interface records it as an
obligation and does not reprove it.

HCP uses this path-equality semantics as motivation for its hidden-
representative model. The public classifier, the cryptographic hiding of a
representative or path witness, and the verifier's evidence boundary are
project-specific additions; they do not follow from Helfer's invariance
theorem alone.

\subsection{Public Classifier Fibers}

Let \(A\) be a type of hidden representatives and let \(P\) be a type of public
class data. A public classifier is a map
\[
  \pi:A\longrightarrow P.
\]
It determines public fibers
\[
  A_\alpha=\pi^{-1}(\alpha)=\{x\in A:\pi(x)=\alpha\},
  \qquad \alpha\in\image(\pi),
\]
and a public equivalence relation
\[
  x\equiv_{\pi}y
  \quad\Longleftrightarrow\quad
  \pi(x)=\pi(y).
\]
The classifier is path-respecting when
\[
  x\simeq_A y\Longrightarrow x\equiv_{\pi}y,
\]
and nonfaithful when some public fiber contains distinct hidden
representatives. A path \(\gamma\in\mathcal P_A(x,y)\) is a hidden witness
that \(x\) and \(y\) lie in the same path component; under path-respect, its
public consequence is the equality \(\pi(x)=\pi(y)\). At this stage the public
datum is only the classifier value. The finite-code projection and
Reed--Solomon/MCA endpoint are introduced later in
\cref{sec:rs-proximity-projection}.

Path-respecting gives only the displayed forward implication. Unless a
separate completeness condition is imposed, equality of classifier values
need not imply the existence of a path; a public fiber \(A_\alpha\) may be a
coarser union of path components. Accordingly, \(\alpha\) denotes the public
classifier value and \(A_\alpha\) its fiber, not a hidden path witness.

Nor do path-respect and nonfaithfulness prove that \(\pi\) hides \(x\):
nonfaithfulness supplies a collision, but it does not imply one-wayness,
indistinguishability, or the absence of other private information from the
output. Such privacy is a separate cryptographic property of the chosen
public-view construction and its distributions.

The fibration projection and public classifier describe compatible but
distinct levels. When the hidden representatives are the objects of \(E\), a
map \(\kappa:B\to P\) on objects supplies the compatibility condition
\[
  \pi=\kappa\circ p.
\]
Then \(p(x)\) is the fibration-base point of \(x\), while
\(\alpha=\pi(x)=\kappa(p(x))\) is its public class label. The special case
\(B=P\) and \(\kappa=\mathrm{id}_B\) identifies the underlying object maps
\(\pi\) and \(p\). In general, \(\kappa\) may be lossy, and
\[
  A_\alpha
  =\bigcup_{\substack{b\in B\\ \kappa(b)=\alpha}}E_b.
\]
This compatibility datum is additional structure; it is not implied merely by
the requirement that \(\pi\) respect paths.

\subsection{Axiom Interfaces}

The hypotheses used below are grouped by mathematical interface. This grouping
does not assert that every item below is literally a file under \lean{Axioms/}.
The relevant interfaces are the following.
\begin{enumerate}[label=(\roman*)]
  \item Foundational cochain/Hodge interface: a cochain complex, a coboundary,
  the induced Hodge--Laplacian, Hodge decomposition, and orthogonality.
  \item Hidden semantic interface: paths, transport, fibration structure,
  homotopy classes, public classifiers, and the nonfaithfulness of public
  class data.
  \item Sheaf-coherence interface: restriction-shape compatibility, naturality,
  \(d^{k+1}d^k=0\), the induced Hodge--Laplacian identity, finiteness,
  simpliciality, and connectedness.
  \item Spectral and stage-admissibility interface: Tsirelson-style spectral
  window conditions, near-boundary conditions, and object-local admissibility
  wrappers for candidate Hodge--Laplacian data.
  \item Public-judgment interface: finite Heyting collapse, public truth
  equivalence, oracle query/answer conditions, and oracle-visible spectral,
  identity, and topology judgments.
  \item Public-artifact interface: cochain-commitment payloads, commitment
  security assumptions, proof-system assumptions, proof-carrying artifacts, and
  compression/observation boundaries used in \cref{sec:public-artifacts-commitments}.
\end{enumerate}
The well-formedness boundaries are displayed below as enumerated conditions.

\begin{definition}[Helfer-derived semantic boundary]\label{def:helfer-semantic-boundary}
  A first-order homotopical semantic datum is admissible if and only if all
  of the following hold:
  \begin{enumerate}[label=(\roman*)]
    \item Equality is interpreted by the path/diagonal condition
      \(x\simeq_A y\Longleftrightarrow\mathcal P_A(x,y)\neq\varnothing\).
    \item Cartesian lifting is available for the relevant fibrations.
    \item Dependent sums and products are coherent with transport and
      base-change along the fibration.
    \item Homotopy-equivalent structures have the same first-order theory.
    \item Class membership is witnessed by hidden path data, and under a
      path-respecting public classifier its public consequence is
      \(\pi(x)=\pi(y)\).
  \end{enumerate}
  Clauses (i)--(iv) extract the portion of Helfer's semantics used here;
  clause (v) is the HCP public-classifier specialization and is not a theorem
  stated in Helfer's paper. The path/classifier checklist corresponds to the
  stable repository interface \lean{Axioms.HelferFOL} and its conjunction
  \lean{Axioms.HelferFOLWellFormed}. The typed interpretation and equivalence
  layer is \lean{Hcrypt.Helfer.Signature},
  \lean{Hcrypt.Helfer.Interpretation}, and
  \lean{Hcrypt.Helfer.HomotopyEquivalent}; the external closed-formula theorem
  is exposed as \lean{Hcrypt.Helfer.ClosedFormulaInvariance}. Together these
  are still not a complete Lean formalization of Helfer's theory.
\end{definition}

\begin{definition}[Foundation assumption boundary]\label{def:foundation-assumption-boundary}
  A foundation datum is admissible if and only if all of the following hold:
  \begin{enumerate}[label=(\roman*)]
    \item \(C^\bullet\) is a cochain complex.
    \item Each coboundary \(d^k:C^k\to C^{k+1}\) is linear.
    \item The Hodge--Laplacian is obtained from adjacent coboundaries and
      adjoints.
    \item \(C^1\) admits the Hodge decomposition
      \[
        C^1=\mathcal H^1\oplus\image(\delta_0)\oplus\image(\delta_1^\ast).
      \]
    \item The summands
      \(\mathcal H^1\), \(\image(\delta_0)\), and
      \(\image(\delta_1^\ast)\) are pairwise orthogonal.
  \end{enumerate}
  These clauses correspond to \lean{Axioms.Foundation} and
  \lean{Axioms.FoundationWellFormed}.
  Commitment authentication, hiding, homomorphism, context binding, and
  partial-opening soundness belong to the separate commitment-security
  boundary below; none is inherited merely from this foundation datum.
\end{definition}

\begin{definition}[Sheaf-coherence assumption boundary]\label{def:sheaf-coherence-assumption-boundary}
  A sheaf-coherence datum is admissible if and only if all of the following
  hold:
  \begin{enumerate}[label=(\roman*)]
    \item Restriction maps have compatible source and target dimensions on
      incidences.
    \item For every edge, the two restriction/transport paths commute.
    \item On every triangle, the face-composition maps agree.
    \item Adjacent sheaf coboundaries square to zero:
      \(d^{k+1}d^k=0\).
    \item The Hodge--Laplacian is obtained from the sheaf coboundaries and
      adjoints.
    \item The harmonic, exact, and coexact summands are pairwise orthogonal.
    \item The nerve is finite.
    \item The nerve is simplicial.
    \item The nerve is connected.
  \end{enumerate}
  These clauses correspond to \lean{Axioms.SheafCoherence} and
  \lean{Axioms.SheafCoherenceWellFormed}.
\end{definition}

The preceding boundary is a condition on the ambient sheaf complex.  A
particular serialized commitment requires an additional, element-dependent
condition.  Let
\[
  \operatorname{dec}_S:\mathcal C_{\mathrm{com}}\longrightarrow C_S^1
\]
be the interpretation of commitment payloads in the degree-one cochain carrier
of \(S\), and let \(c\mathrel{\mathsf{binds}}S\) denote the topology-binding
condition for that interpretation.

\begin{definition}[Committed one-cochain coherence]\label{def:committed-one-cochain-coherence}
  A commitment payload \(c\) is coherent with a sheaf datum \(S\), written
  \(c\in\mathcal G_S\), if and only if
  \begin{enumerate}[label=(\roman*)]
    \item \(c\) is well formed as a serialized commitment payload;
    \item \(S\) satisfies the sheaf-coherence boundary of
      \cref{def:sheaf-coherence-assumption-boundary};
    \item \(c\mathrel{\mathsf{binds}}S\); and
    \item there is an \(a\in C_S^0\) such that
      \[
        \operatorname{dec}_S(c)=\delta_0 a.
      \]
  \end{enumerate}
  Thus \(c\) represents an exact one-cochain, rather than merely having a
  valid wire shape.
\end{definition}

\begin{proposition}[Committed exact cochains are cocycles]
\label{prop:committed-exact-cochains-are-cocycles}
  If \(c\in\mathcal G_S\), then
  \[
    \delta_1\operatorname{dec}_S(c)=0.
  \]
\end{proposition}

\begin{proof}
  Choose \(a\in C_S^0\) with
  \(\operatorname{dec}_S(c)=\delta_0a\).  Ambient sheaf coherence gives the
  actual cochain-map equation \(\delta_1\delta_0=0\).  Hence
  \[
    \delta_1\operatorname{dec}_S(c)
    =\delta_1\delta_0a
    =0.
  \]
\end{proof}

\begin{definition}[Proof-system and commitment-security boundaries]\label{def:proof-commitment-assumption-boundary}
  The commitment-security boundary holds if and only if all of the following
  hold:
  \begin{enumerate}[label=(\roman*)]
    \item Exact-payload binding: the authenticated payload binds the serialized
      \(C^1\) value to its topology/instance identifier. If two degree-zero
      openings produce that value, their difference lies in \(\ker\delta_0\);
      this is the intended opening equivalence, not a uniqueness theorem.
    \item Hiding: the commitment distribution hides the committed witness
      component.
    \item Homomorphic addition:
      \(C(m_1,r_1)+C(m_2,r_2)=C(m_1+m_2,r_1+r_2)\) when the two commitments
      are compatible.
    \item Homomorphic scalar multiplication:
      \(a\,C(m,r)=C(am,ar)\) for admissible scalars \(a\).
    \item Topology binding: a commitment is bound to its underlying
      simplicial complex or topology identifier.
    \item Partial opening soundness: opening a specified subset verifies the
      corresponding local cochain constraints.
  \end{enumerate}
  The proof-system boundary holds if and only if all of the following hold:
  \begin{enumerate}[label=(\roman*)]
    \item Completeness: an honest proof for a valid witness verifies.
    \item Soundness: verifier acceptance implies validity of the asserted
      mathematical statement, up to the stated error model.
    \item Zero knowledge: the proof transcript does not identify the hidden
      witness.
    \item Statement binding: the proof is bound to its public statement,
      topology, and instance context.
  \end{enumerate}
  These clauses correspond to
  \lean{Axioms.CochainCommitment.CochainCommitmentAxioms},
  \lean{Axioms.CochainCommitment.CochainCommitmentWellFormed},
  \lean{Axioms.ZKP}, and \lean{Axioms.ZKPWellFormed}.
\end{definition}

Consensus, ledger, reliability, and Proof-of-Value records are mentioned here
only at their mathematical interface boundary. Their runtime policy content
belongs to the subsequent protocol paper.

\subsection{Exact Structural Invariants}

The exact structural invariants are not a separate source of geometry. They
package the same foundational and sheaf-coherence hypotheses at the point
where an object is admitted. We use the cellular-sheaf Hodge--Laplacian
viewpoint of Hansen--Ghrist as background \cite{hansenGhrist2018sheaves}. For
a finite cell poset \(P_X\), ordered by the face relation, a cellular sheaf is
represented here in the face-to-coface convention as
\[
  P_X\xrightarrow{F}\mathsf{Vect}_{\R},
  \qquad
  C^0(X;F)\xrightarrow{\delta_0}C^1(X;F)
  \xrightarrow{\delta_1}C^2(X;F).
\]
For a cochain complex \(C^\bullet\), the accepted data supply differentials
\[
  d^k:C^k\longrightarrow C^{k+1}
\]
with \(d^{k+1}d^k=0\). Together with adjoints, these give the degree-\(k\)
Hodge--Laplacian
\[
  \Delta_k=d^{k-1}(d^{k-1})^\ast+(d^k)^\ast d^k.
\]
In the sheaf setting, the corresponding coherence assumptions assert that
restriction maps, local data, and gluing data are compatible with this
cochain/Hodge structure. The only structural consequences used later are the
adjacent differential identity \(d^{k+1}d^k=0\) and the Hodge--Laplacian
identity above.
For the projection chain, we additionally require the maps used to pass from
hidden sheaf data to public cochain-facing data to commute with the
coboundary:
\[
\begin{array}{ccc}
  F(V) & \xrightarrow{\eta_V} & G(V)\\
  \delta_F\downarrow\phantom{\delta_F} & & \phantom{\delta_G}\downarrow\delta_G\\
  F(E) & \xrightarrow{\eta_E} & G(E)
\end{array}
\qquad
\eta_E\circ\delta_F=\delta_G\circ\eta_V .
\]
This is the functorial/naturality diagram used in the paper. It is a
cochain-level coherence condition, not a runtime statement about byte transport
or consensus execution.

\begin{proposition}[Exact structural preconditions]\label{prop:exact-structural-preconditions}
  If a sheaf exact-gate bundle is accepted as well formed, then its
  sheaf-coherence data supply
  \[
    d^{k+1}d^k=0
    \qquad\text{and}\qquad
    \Delta_k=d^{k-1}(d^{k-1})^\ast+(d^k)^\ast d^k
  \]
  for the cochain complex under discussion.
\end{proposition}

\begin{proof}
  Well-formedness of the exact-gate bundle includes the sheaf-coherence
  condition. By definition of that condition, the adjacent coboundaries
  compose to zero and the recorded Laplacian is the displayed sum of the two
  coboundary terms. Projecting these two fields from the well-formedness
  witness gives the stated identities.
\end{proof}

\subsection{Spectral-Gap Admissibility}

The spectral condition is an object-level admissibility condition. It is not the
source of HCP hardness, not a reconstruction theorem for the hidden fibration,
and not introduced here as a consensus rule. In this paper, the
Proof-of-Value (PoV)/Tsirelson
spectral window is a mathematical condition on candidate
degree-one Hodge--Laplacian data; its use as a protocol-level admissibility rule belongs to
the subsequent protocol paper. The terminology is
Tsirelson-style because the admissible window uses the familiar \(\sqrt 2\)
scale from Tsirelson's bound \cite{tsirelson1980bell}. For a Hermitian
positive semidefinite operator \(\Delta\), write
\[
  \lambda_1^+(\Delta)
  =
  \inf\{\lambda:\lambda\text{ is an eigenvalue of }\Delta,\ \lambda\ne 0\},
\]
with \(\lambda_1^+(\Delta)=0\) when \(\Delta\) has no nonzero spectral value.
The Tsirelson-style admissibility axiom is the spectral window
\[
  \frac{\sqrt{2}}{2}
  \le
  \lambda_1^+(\Delta)
  \le
  2\sqrt{2}.
\]
A separate boundary condition marks \(\Delta\) as \(\varepsilon\)-near one of
the two endpoints:
\[
  \left|\lambda_1^+(\Delta)-\frac{\sqrt 2}{2}\right|\le \varepsilon
  \quad\text{or}\quad
  \left|\lambda_1^+(\Delta)-2\sqrt 2\right|\le \varepsilon .
\]
This condition records near-boundary status only; runtime policy for such values
belongs to the subsequent protocol paper.

The corresponding Grassmannian/Laplacian interface makes the hidden-kernel
search space concrete. Throughout this subsection, set \(R=\mathbb R\), with its
standard inner product. Let
\[
  \mathcal L_\beta(R^n)
  =
  \{\Delta:\Delta=\Delta^\ast,\ \Delta\succeq 0,\ \dim_R\ker\Delta=\beta\}
\]
be the Laplacian space at kernel dimension \(\beta\). For
\(W\in\Gr(k,R^n)\), let \(P_W\) be the orthogonal projection onto \(W\) and set
\[
  \Delta_W = I-P_W .
\]
Thus \(\ker(\Delta_W)=W\), and the Grassmannian family is
\[
  \Omega_k
  =
  \{\Delta_W:W\in\Gr(k,R^n)\}.
\]
Equivalently, a point of \(\Gr(k,R^n)\) is a \(k\)-dimensional harmonic
subspace \(W\), and the map \(W\mapsto\Delta_W\) sends it to the corresponding
projection Laplacian. In particular, \(\Omega_1\) is the Grassmannian family
of one-dimensional harmonic kernels.

\begin{proposition}[Grassmannian Laplacian search space]\label{prop:grassmannian-search-space}
  For every \(k\le n\), the family \(\Omega_k\) is nonempty and satisfies
  \[
    \Omega_k
    \subseteq
    \{\Delta_W:W\subset R^n\text{ a linear subspace}\}.
  \]
  Moreover, every element of \(\Omega_k\) has a \(k\)-dimensional harmonic
  kernel, so \(\Omega_k\subseteq\mathcal L_k(R^n)\).
\end{proposition}

\begin{proof}
  If \(k\le n\), choose \(k\) coordinate basis vectors in \(R^n\). Their span
  is a \(k\)-dimensional subspace \(W\), and \(\Delta_W\in\Omega_k\), proving
  nonemptiness. The displayed containment is immediate from the definition of
  \(\Omega_k\). Since \(\Delta_W=I-P_W=P_{W^\perp}\), its kernel is exactly
  \(W\), hence has dimension \(k\).
\end{proof}

\begin{proposition}[Injectivity of the Grassmannian Laplacian map]\label{prop:grassmannian-laplacian-injective}
  If \(W_1,W_2\subset R^n\) are subspaces and
  \[
    \Delta_{W_1}=\Delta_{W_2},
  \]
  then \(W_1=W_2\). Consequently \(W\mapsto\Delta_W\) is injective on
  \(\Gr(k,R^n)\).
\end{proposition}

\begin{proof}
  The equality \(\Delta_{W_1}=\Delta_{W_2}\) is the equality
  \(I-P_{W_1}=I-P_{W_2}\), hence \(P_{W_1}=P_{W_2}\). For any
  \(x\in W_1\), \(P_{W_1}x=x\), so \(P_{W_2}x=x\), and therefore
  \(x\in W_2\). Thus \(W_1\subseteq W_2\). The reverse containment is
  identical, so \(W_1=W_2\).
\end{proof}

The higher-rank Grassmannian search spaces contain explicit one-parameter
families. Fix \(1\le k<n\). Choose linearly independent vectors \(u,v\in R^n\)
and a \((k-1)\)-dimensional subspace \(W_0\) disjoint from
\(\operatorname{span}\{u,v\}\). For \(t\in R\), set
\[
  W(t)=\operatorname{span}\{u+t v\}+W_0.
\]
Then \(W(t)\) has dimension \(k\), and varying \(t\) gives a corresponding
family
\[
  \Delta_{W(t)}=I-P_{W(t)}.
\]

\begin{proposition}[Parametric Laplacian search space]\label{prop:parametric-laplacian-search-space}
  If \(1\le k<n\), then there exist \(u,v,W_0\) as above such that the map
  \[
    R\longrightarrow\Gr(k,R^n),
    \qquad
    t\longmapsto W(t)
  \]
  is injective. Hence \(t\mapsto\Delta_{W(t)}\) is an injective
  one-parameter family in \(\Omega_k\).
\end{proposition}

\begin{proof}
  The existence of \(u,v,W_0\) follows by choosing a \((k-1)\)-plane and then
  two independent directions outside it; the inequality \(k<n\) leaves enough
  room for this choice. The dimension computation uses the disjointness:
  \[
    \dim W(t)
    =
    \dim\operatorname{span}\{u+tv\}+\dim W_0
    =
    1+(k-1)=k.
  \]
  If \(W(t_1)=W(t_2)\), then \(u+t_1v\) lies in
  \(\operatorname{span}\{u+t_2v\}+W_0\). Modulo \(W_0\), this gives a linear
  dependence between \(u\) and \(v\) unless \(t_1=t_2\). Since \(u,v\) were
  chosen independent and \(W_0\) was disjoint from their span, \(t_1=t_2\).
  Injectivity of \(t\mapsto\Delta_{W(t)}\) follows from
  \cref{prop:grassmannian-laplacian-injective}.
\end{proof}

\begin{proposition}[Projection spectrum and spectral gap]\label{prop:projection-spectrum-gap}
  For every \(W\subset R^n\), the projection Laplacian
  \(\Delta_W=I-P_W=P_{W^\perp}\) has spectrum contained in \(\{0,1\}\).
  Its kernel is \(W\), its image is \(W^\perp\), and, if
  \(W\ne R^n\), its least nonzero eigenvalue and largest eigenvalue are both
  \(1\). Thus the relative spectral gap is \(1\) for the Grassmannian
  projection family with \(k<n\).
\end{proposition}

\begin{proof}
  If \(x\in W\), then \(\Delta_Wx=0\). If \(x\in W^\perp\), then
  \(\Delta_Wx=x\). Since \(R^n=W\oplus W^\perp\), these two subspaces exhaust
  the eigenspace decomposition. Therefore the only possible eigenvalues are
  \(0\) and \(1\), with \(0\)-eigenspace \(W\) and \(1\)-eigenspace
  \(W^\perp\). When \(W^\perp\ne 0\), the nonzero spectral values consist
  only of \(1\), so both \(\lambda_1^+\) and \(\lambda_{\max}\) are \(1\).
\end{proof}

The likelihood surface records the corresponding adversary vocabulary. A
classical adversary is a map from an observation \(b\) to a guessed Laplacian;
the density-based interface assigns a probability density to candidate
Laplacians. For a
noise scale \(\sigma>0\), the Gaussian likelihood of a candidate \(A\) at
observation \(b\) is written abstractly as
\[
  L_\sigma(b,A)
  =
  \exp\left(-\frac{\|A b\|_2^2}{2\sigma^2}\right),
\]
with normalized classical and density-based likelihood scores obtained by
integrating or maximizing this likelihood over the chosen family
\(\Omega_k\). These likelihood objects are comparison surfaces for quantitative
security statements; they are not assumptions used to prove the hidden-operator
separation in \cref{prop:cvp-insufficient}.

\begin{definition}[Adversary likelihood surfaces]\label{def:adversary-likelihood-surfaces}
  A classical adversary on \(R^n\) is a function
  \[
    \mathcal A_{\mathrm{cl}}:R^n\to M_n(R)
  \]
  from an observation to a guessed Laplacian. The interface called a
  \emph{quantum adversary} in the Lean development is a nonnegative density
  on a candidate family \(\Omega_k\):
  \[
    \mathcal A_{\mathrm q}:R^n\times\Omega_k\to R_{\ge 0}.
  \]
  Its integral over \(\Omega_k\), with respect to a chosen measure \(\mu\),
  is \(1\) for each observation. With the standard Gaussian prior on
  observations, the corresponding normalized likelihood scores are obtained
  from
  \[
    \int_{R^n}
      L_\sigma\bigl(b,\mathcal A_{\mathrm{cl}}(b)\bigr)
      e^{-\|b\|_2^2/2}\,db
  \]
  in the classical case and
  \[
    \int_{R^n}
      \left(
        \int_{\Omega_k}
          \mathcal A_{\mathrm q}(b,A)L_\sigma(b,A)\,d\mu(A)
      \right)
      e^{-\|b\|_2^2/2}\,db
  \]
  in the quantum-density case, up to the usual Gaussian normalization factor.
\end{definition}

These scores measure consistency with an observation. They do not by
themselves measure recovery of the operative hidden operator or formalize a
quantum computation model.

The likelihood is written with the Euclidean norm. The norm-convention check is
that the default function-space norm may be the supremum norm, so the Euclidean
likelihood must be represented by passing to the \(\ell^2\) norm explicitly. The
mathematical surface used here is the corrected Euclidean one.
The corrected pointwise likelihood and Gaussian factorization have Lean
endpoints. The displayed Euclidean integrated scores are not identified with
the legacy \lean{success_prob_*} functions, which use the function-space
supremum-norm likelihood; those integrated endpoints are outside the selected
Hodge/CVP entry.

\begin{proposition}[Likelihood for a projection Laplacian]\label{prop:likelihood-projection-laplacian}
  For \(A=\Delta_W\), the Euclidean Gaussian likelihood is
  \[
    L_\sigma(b,\Delta_W)
    =
    \exp\left(
      -\frac{\|P_{W^\perp}b\|_2^2}{2\sigma^2}
    \right).
  \]
\end{proposition}

\begin{proof}
  By \cref{prop:projection-harmonicity-equivalences},
  \(\Delta_Wb=P_{W^\perp}b\). Substitute this identity into the definition of
  \(L_\sigma(b,A)\).
\end{proof}

\begin{proposition}[Single-observation optimal classical guess]\label{prop:single-observation-optimal-guess}
  If \(b\ne 0\) and \(W_b=\operatorname{span}\{b\}\), then
  \[
    L_\sigma(b,\Delta_{W_b})=1.
  \]
  Equivalently, the projection residual of \(b\) against the Laplacian whose
  kernel is \(\operatorname{span}\{b\}\) is zero.
\end{proposition}

\begin{proof}
  Since \(b\in W_b\), \(\Delta_{W_b}b=0\). Hence
  \(L_\sigma(b,\Delta_{W_b})=\exp(0)=1\).
\end{proof}

This candidate achieves maximum consistency with each observation, but it
need not be the operative hidden kernel. Likelihood maximization alone
therefore does not recover the HCP target.

\begin{proposition}[Gaussian prior factorization]\label{prop:gaussian-prior-factorization}
  For \(x=(x_i)_{i=1}^n\in R^n\),
  \[
    \exp\left(-\frac{\|x\|_2^2}{2}\right)
    =
    \prod_{i=1}^n
    \exp\left(-\frac{x_i^2}{2}\right).
  \]
\end{proposition}

\begin{proof}
  Since \(\|x\|_2^2=\sum_i x_i^2\), the left hand side is
  \[
    \exp\left(\sum_i -\frac{x_i^2}{2}\right),
  \]
  which equals the product of the exponentials.
\end{proof}

\begin{definition}[Spectral fingerprint]\label{def:spectral-fingerprint}
  Let \(L:V\to V\) be a self-adjoint positive semidefinite operator on a
  finite-dimensional inner-product space, such as a cellular sheaf Laplacian.
  Its spectral fingerprint is the multiset
  \[
    \operatorname{Spec}_{\mathrm{multi}}(L)
  \]
  of eigenvalues counted with multiplicity. Two Laplacians are spectrally
  equivalent when these multisets agree. A separate kernel-rank relation
  records only harmonic dimension:
  \[
    L_1\equiv_{\dim\ker}L_2
    \quad\Longleftrightarrow\quad
    \dim\ker L_1=\dim\ker L_2.
  \]
  This relation is not Helfer homotopy equivalence and is not a chain-isometric
  equivalence of hidden presentations.
\end{definition}

The order of determination is important. The public classifier value
\(\alpha\) gives the public fiber \(A_\alpha\); it is not itself a fibration
ID. A hidden representative \(x\in A_\alpha\) determines the operative
Hodge--Laplacian \(\Delta_x\). Once an HCP witness has selected or certified
the hidden Hodge data \((\Delta_x,h,e)\), write
\[
  \operatorname{FID}(x):=
  \operatorname{Spec}_{\mathrm{multi}}(\Delta_x).
\]
This is a spectral-class tag in the certified Hodge layer, not an injective
identifier of \(x\) or of its full presentation. In particular, isometrically
conjugate Laplacians have the same fingerprint. The tag remains downstream of
harmonic-cochain recovery and is not part of the public classifier value.

\begin{proposition}[Spectrum is not determined by harmonic dimension]\label{prop:spectrum-not-determined-by-kernel-dimension}
  There exist self-adjoint positive semidefinite Laplacians \(L_1,L_2\) with
  \[
    \dim\ker L_1=\dim\ker L_2
  \]
  but
  \[
    \operatorname{Spec}_{\mathrm{multi}}(L_1)
    \ne
    \operatorname{Spec}_{\mathrm{multi}}(L_2).
  \]
\end{proposition}

\begin{proof}
  Take \(V=R\), \(L_1=\operatorname{id}_R\), and
  \(L_2=2\operatorname{id}_R\). Both kernels are zero, so the harmonic
  dimensions agree. The spectral fingerprints are the singleton multisets
  \(\{1\}\) and \(\{2\}\), respectively, so the operators are not spectrally
  equivalent.
\end{proof}

\begin{proposition}[Spectral-gap admissibility]\label{prop:spectral-admissibility}
  If a candidate Hodge--Laplacian geometry carries a by-construction
  PoV/Tsirelson admissibility witness, then
  \[
    \lambda_1^+(\Delta)\in
    \left[\frac{\sqrt{2}}{2},\,2\sqrt{2}\right].
  \]
\end{proposition}

\begin{proof}
  The by-construction witness contains
  \[
    \frac{\sqrt{2}}{2}\le\lambda_1^+(\Delta)
    \qquad\text{and}\qquad
    \lambda_1^+(\Delta)\le 2\sqrt{2}.
  \]
  These two inequalities are exactly membership in the displayed closed
  interval.
\end{proof}

Thus the by-construction witness is an admissibility certificate for the
displayed spectral window. It constrains which candidate objects are admissible;
it does not remove hidden-operator ambiguity. The HCP hardness mechanism in
\cref{sec:hcp-definition} remains the hidden fibration, hidden Laplacian, and
hidden kernel.

For a nonempty family \(\Omega\subseteq\mathcal L_1(R^n)\), the associated
classical and quantum adversary-bound forms are spectral-statistical targets,
not HCP hardness assumptions. The corresponding empty-family statements are
only boundary checks; they are not nontrivial hardness bounds.

\subsection{Heyting and Oracle Judgment Surface}

Let \(\Omega\) be the carrier of hidden truth values. The ownership split is:
\[
\begin{array}{rcl}
\text{hidden equality} &:&
  x\simeq y\Longleftrightarrow \mathcal P(x,y)\neq\varnothing,\\[2pt]
\text{public class data} &:&
  x\equiv_\pi y\Longleftrightarrow \pi(x)=\pi(y),\\[2pt]
\text{public truth data} &:&
  q:\Omega\longrightarrow\Omega_3 .
\end{array}
\]
The first line is the Helfer path/diagonal interface
\cite[Section~2]{helfer2019fol}. The second line is the public classifier interface
defined above. The third line is the finite public judgment interface. We take
\[
  \Omega_3=\{\bot,\, ?,\, \top\},
  \qquad
  \bot< ? <\top,
\]
with meet and join given by minimum and maximum in the chain. Its implication
is the residuated operation characterized by
\[
  c\le (a\Rightarrow b)
  \quad\Longleftrightarrow\quad
  a\wedge c\le b.
\]
We assume that the public collapse \(q\) is nonfaithful, in the sense that
\[
  \exists\,\omega_1\ne\omega_2
  \quad\text{with}\quad
  q(\omega_1)=q(\omega_2).
\]

The oracle-specific part is the query/answer mechanism that produces these
public judgments. For this boundary we use oracle-modality terminology
compatible with Ahman--Bauer \cite{ahmanBauer2026oracle}: oracle-visible
conditions such as vertex compatibility, Betti-number compatibility, public
identity-hash compatibility, and spectral compatibility are evaluated in the
finite public Heyting algebra.
Thus an oracle answer is a public quotient of hidden semantic evidence, not a
reconstruction of the hidden representative, hidden path witness, fibration,
Laplacian, or harmonic kernel.

Operational maps out of \(\Omega_3\), such as protocol actions attached to the
intermediate public judgment \(?\), belong to the subsequent protocol paper's
runtime and consensus semantics. They are not the Heyting implication operation
and do not alter the residuation law above.

The identity hash used by the oracle interface is an observation obtained from
the public Hodge projection and quantization path. It is not a cochain
commitment payload. The corresponding stability assumption says that equal
public identity inputs give the same identity hash; it is a determinism
condition for public identity observations, not a uniqueness or binding theorem
for cochain commitments.
Similarly, topological observations such as Betti-number compatibility are
homotopy-invariant public judgments. Spectral fingerprints are post-recovery
spectral-class tags: they are not determined by harmonic dimension alone by
\cref{prop:spectrum-not-determined-by-kernel-dimension}. The oracle interface
therefore keeps Betti/topology observations, spectral admissibility, identity
hashes, and cochain commitments as separate public conditions.

This modular split is part of the mathematical statement: cryptographic
assumptions, public classifiers, and Hodge-theoretic claims are not identified.
Exact and oracle invariants are used here as mathematical acceptance conditions;
their full computational and protocol treatment is deferred to the subsequent
protocol paper.

\subsection*{Lean Cross-Reference}

\begin{leanledger}
The semantic and axiom interfaces in this section correspond to
\lean{Sheaf.OrientedTwoComplex}, \lean{Sheaf.CellularSheaf},
\lean{Sheaf.CellularSheaf.Morphism},
\lean{Sheaf.CellularSheaf.Morphism.mapC1_delta0},
\lean{Sheaf.CellularSheaf.Morphism.mapC2_delta1},
\lean{Sheaf.CellularSheaf.delta1_comp_delta0},
\lean{Sheaf.CellularSheaf.degreeOneLaplacian},
\lean{Sheaf.CellularSheaf.ker_degreeOneLaplacian},
\lean{Axioms.HelferFOL}, \lean{Axioms.Fibration},
\lean{Hcrypt.Helfer.Signature}, \lean{Hcrypt.Helfer.Interpretation},
\lean{Hcrypt.Helfer.HomotopyHomomorphism},
\lean{Hcrypt.Helfer.HomotopyEquivalent},
\lean{Hcrypt.Helfer.ClosedFormulaInvariance},
\lean{Axioms.PublicClassifierFactorization},
\lean{Axioms.PublicClassifierFactorization.publicFiber_eq_iUnion_baseFibers},
\lean{Axioms.HomotopyClass}, \lean{Axioms.Foundation},
\lean{Axioms.FoundationWellFormed}, \lean{Axioms.SheafCoherence},
\lean{Axioms.SheafCoherenceWellFormed},
\lean{Axioms.CommittedCochainWellFormed},
\lean{Axioms.RepresentsExactOneCochain},
\lean{Axioms.IsOneCocycle},
\lean{Axioms.oneCocycle_of_committedCochainWellFormed}, \lean{Gates.Exact},
\lean{Gates.structural_preconditions_of_sheaf_exact_gates},
\lean{Axioms.SpectralBounds}, \lean{Tsirelson},
\lean{Gates.Spectral}, \lean{Axioms.StageAdmissibility},
\lean{Gates.Stage}, \lean{Axioms.Heyting}, and
\lean{Axioms.OracleModality}.
Here \lean{Axioms.HelferFOL} is an extracted specification boundary: Lean
proves the local reflexivity, symmetry, transitivity, fiber-membership, and
classifier consequences built from its path interface. The typed Helfer
modules formalize the algebraic-interpretation equivalence used by HCP, but
Helfer's general formula-in-context invariance proof remains external; its
closed-formula consequence is passed explicitly to the Lean bridge rather than
asserted globally.
\end{leanledger}
